\begingroup
\renewcommand{\arraystretch}{1.45}
\begingroup
\Needspace{45\baselineskip}
\InterpolationArticle{31}
Since the rules explained above provide a function $X'$ of the form
\[
\begin{aligned}
&\gamma+\gamma'\cos x+\gamma''\cos2x+.\,.\,.+\gamma^{m}\cos mx\\[5pt]
&\quad+\delta'\sin x+\delta''\sin2x+.\,.\,.\,.+\delta^{m}\sin mx
\end{aligned}
\]
satisfying all $\mu$ given values, and in which $m=\frac12\mu-\frac12$ or $\frac12\mu$ according as $\mu$ is odd or even, it is worthwhile to compare this function with the function just found
\[
\begin{aligned}
&\varepsilon+\varepsilon'\cos x+\varepsilon''\cos2x+\varepsilon'''\cos3x+.\,.\,.\\[5pt]
&\qquad{}+\varepsilon^{\mu-1}\cos(\mu-1)x
\end{aligned}
\]
\endgroup
\SourceLeaf{322}
which we shall denote by $X''$. Comparing the values of the coefficients
$\gamma,\gamma',\delta',\gamma'',\delta''$ etc. given in Articles 20, 22 with the values of the coefficients
$\varepsilon,\varepsilon',\varepsilon''$ etc. just found gives
\[
\begin{aligned}
&\varepsilon=\gamma\\[5pt]
&\varepsilon'=\gamma'-\delta'\operatorname{cotang}\mu a,
\\[5pt]
&\varepsilon^{\mu-1}=\delta'\operatorname{cosec}\mu a\\[5pt]
&\varepsilon''=\gamma''-\delta''\operatorname{cotang}\mu a,
\\[5pt]
&\varepsilon^{\mu-2}=\delta''\operatorname{cosec}\mu a\\[5pt]
&\varepsilon'''=\gamma'''-\delta'''\operatorname{cotang}\mu a,
\\[5pt]
&\varepsilon^{\mu-3}=\delta'''\operatorname{cosec}\mu a\\[5pt]
&\text{etc., as far as}\\[5pt]
&\left(\tfrac12\right)\varepsilon^m=\gamma^m-\delta^m\operatorname{cotang}\mu a,
\\[5pt]
&\left(\tfrac12\right)\varepsilon^{\mu-m}=\delta^m\operatorname{cosec}\mu a
\end{aligned}
\]
\Needspace{20\baselineskip}
The factor $\tfrac12$ placed before $\varepsilon^m$ and $\varepsilon^{\mu-m}$ applies only when $\mu$ is even and must be omitted when $\mu$ is odd. In this way, in the latter case, where $\mu-m=m+1$, all $\mu$ coefficients $\varepsilon,\varepsilon',\varepsilon''$ etc. are evidently determined by $\gamma,\gamma',\delta'$ etc.; in the former, however, where $m=\mu-m$, for the coefficient $\varepsilon^{\frac12\mu}$ we have the two values $2\gamma^{\frac12\mu}-2\delta^{\frac12\mu}\cot\mu a$, $2\delta^{\frac12\mu}\operatorname{cosec}\mu a$, whose equality is easily seen from the equation $\gamma^{\frac12\mu}\sin\tfrac12\mu a=\delta^{\frac12\mu}\cos\tfrac12\mu a$, and in place of which we can therefore also adopt this value: $\gamma^{\frac12\mu}-\delta^{\frac12\mu}\operatorname{cotg}\mu a+\delta^{\frac12\mu}\operatorname{cosec}\mu a$.

\Needspace{25\baselineskip}
It is easily deduced from this that, in both cases,
\[
\begin{aligned}
&X''-(X'-\delta'\sin x-\delta''\sin2x-\delta^m\sin3x-\text{etc.}-\delta^m\sin mx)\\[5pt]
&=-\operatorname{cotg}\mu a(\delta'\cos x+\delta''\cos2x+\delta'''\cos3x+\mathinner{.\,.\,.}+\delta^m\cos mx)\\[5pt]
&\quad+\operatorname{cosec}\mu a(\delta'\cos(\mu-1)x+\delta''\cos(\mu-2)x+\mathinner{.\,.\,.}+\delta^m\cos(\mu-m)x)
\end{aligned}
\]
It follows that $X''$ is obtained from $X'$ by substituting for each term $\delta^\lambda\sin\lambda x$
\[
\begin{aligned}
&-\operatorname{cotg}\mu a\,\delta^\lambda\cos\lambda x+\operatorname{cosec}\mu a\,\delta^\lambda\cos(\mu-\lambda)x
\end{aligned}
\]
or
\[
\begin{aligned}
&\frac{-\delta^\lambda\cos\mu a\cos\lambda x+\delta^\lambda\cos(\mu-\lambda)x}{\sin\mu a}
\end{aligned}
\]
We leave to the readers the comparison of this rule with what was explained in Article 24.

\Needspace{28\baselineskip}
\InterpolationArticle{32}
If the function $X$ does not stop at $\cos nx$ but continues further, the following terms being expressed by $\alpha^\mu\cos\mu x+\alpha^{\mu+1}\cos(\mu+1)x+\text{etc.}$, we shall have
\SourceLeaf{323}
\[
\begin{aligned}
&\varepsilon=\alpha+\alpha^\mu\cos\mu a+\alpha^{2\mu}\cos2\mu a+\text{etc.}\\[5pt]
&\varepsilon'=\frac{1}{\sin\mu a}\{\alpha'\sin\mu a+\alpha^{\mu+1}\sin2\mu a+\alpha^{2\mu+1}\sin3\mu a+\text{etc.}\\[5pt]
&\qquad-\alpha^{2\mu-1}\sin\mu a-\alpha^{3\mu-1}\sin2\mu a-\alpha^{4\mu-1}\sin3\mu a-\text{etc.}\}\\[5pt]
&\varepsilon''=\frac{1}{\sin\mu a}\{\alpha''\sin\mu a+\alpha^{\mu+2}\sin2\mu a+\alpha^{2\mu+2}\sin3\mu+\text{etc.}\\[5pt]
&\qquad-\alpha^{2\mu-2}\sin\mu a-\alpha^{3\mu-2}\sin2\mu a-\alpha^{4\mu-2}\sin3\mu a-\text{etc.}\}\\[5pt]
&\text{etc.}\\[5pt]
&\varepsilon^{\rho-1}=\frac{1}{\sin\mu a}\{\alpha^{\mu-1}\sin\mu a+\alpha^{2\mu-1}\sin2\mu a+\alpha^{3\mu-1}\sin3\mu a+\text{etc.}\\[5pt]
&\qquad-\alpha^{\mu+1}\sin\mu a-\alpha^{2\mu+1}\sin2\mu a-\alpha^{3\mu+1}\sin3\mu a-\text{etc.}\}
\end{aligned}
\]
\Needspace{46\baselineskip}
From these formulas one can judge how far the difference of the coefficients $\varepsilon,\varepsilon',\varepsilon''$ etc. from their true values $\alpha,\alpha',\alpha''$ etc. may be neglected. In this case, therefore, the function $X''$ will not be $X$ itself, although it agrees with it in satisfying all $\mu$ given values. All similar functions satisfying these values will, however, be contained in the form $X''+Y(\cos\mu x-\cos\mu a)$, where $Y$ denotes an undetermined function of the arc $x$ of the same form as $X$, namely free of sines. The function $X''$ will be the only one of the $\mu-1^{\mathrm{th}}$ order: every other similar function satisfying the given values will rise to a higher order. If such a function $X'''$ were known from elsewhere, all similar functions representing the given values would also be contained in the form $X'''+Y(\cos\mu x-\cos\mu a)$, and $Y$ could be determined so as to produce a function reduced to the $\mu-1^{\mathrm{th}}$ order, which will be $X''$ itself. The easiest method for this purpose seems to be first to derive from $X'''$ the function $X'$ by Article 24, and from this $X''$ by Article 31. Thus, let any term in $X'''$ be $L\cos\lambda x$, and set $\lambda=k\mu+\lambda'$, so that $\lambda'<\mu$ and $k$ is an integer. Then, by Article 24, for that term there must be written in $X'$
\[
\begin{aligned}
&L\cos k\mu a\cos\lambda'x-L\sin k\mu a\sin\lambda'x,
\\[5pt]
&\qquad\text{if }\lambda'<\tfrac12\mu,\quad\text{but not }\lambda'=0\\[5pt]
&L\cos(k+1)\mu a\cos(\mu-\lambda')x+L\sin(k+1)\mu a\sin(\mu-\lambda')x.
\\[5pt]
&\qquad\text{if }\lambda'>\tfrac12\mu\\[5pt]
&L\cos\tfrac12\mu a\cos(k+\tfrac12)\mu a\cos\tfrac12\mu x\\[5pt]
&\qquad{}+L\sin\tfrac12\mu a\cos(k+\tfrac12)\mu a\sin\tfrac12\mu x,
\\[5pt]
&\qquad\text{if }\lambda'=\tfrac12\mu
\end{aligned}
\]
\Needspace{38\baselineskip}
From this there results in $X''$, by the preceding article, in the first case
\[
\begin{aligned}
&L\cos k\mu a\cos\lambda'x+L\sin k\mu a\operatorname{cotang}\mu a\cos\lambda'x\\[5pt]
&\qquad{}-L\sin k\mu a\operatorname{cosec}\mu a\cos(\mu-\lambda')x
\end{aligned}
\]
in the second case
\[
\begin{aligned}
&L\cos(k+1)\mu a\cos(\mu-\lambda')x-L\sin(k+1)\mu a\operatorname{cotg}\mu a\cos(\mu-\lambda')x\\[5pt]
&\qquad+L\sin(k+1)\mu a\operatorname{cosec}\mu a\cos\lambda'x
\end{aligned}
\]
\SourceLeaf{324}
in the third case
\[
\begin{aligned}
&L\cos\tfrac12\mu a\cos(k+\tfrac12)\mu a\cos\tfrac12\mu x\\[5pt]
&\qquad{}-L\sin\tfrac12\mu a\cos(k+\tfrac12)\mu a\operatorname{cotg}\mu a\cos\tfrac12\mu x\\[5pt]
&\qquad+L\sin\tfrac12\mu a\cos(k+\tfrac12)\mu a\operatorname{cosec}\mu a\cos\tfrac12\mu x
\end{aligned}
\]
which expression in all three cases reduces to
\[
\begin{aligned}
&\frac{L\sin(k+1)\mu a\cos\lambda'x-L\sin k\mu a\cos(\mu-\lambda')x}{\sin\mu a}
\end{aligned}
\]
Whenever $\lambda'=0$, however, we have in $X'$ simply $L\cos k\mu a$, which term is retained unchanged in $X''$.

\endgroup
